% Source: https://artofproblemsolving.com/community/c3h1774078_easiest_proof_ever
% Stable identity: aops-topic-1774078-q01
We use Bertrand's postulate: for every integer $n\ge2$ there is a prime $p$ with $\dfrac {n}{2}<p\le n$ (for $n=2,3$ take $p=2,3$).

Such a prime divides $n!$, because $p\le n$. It divides $n!$ only once: the multiples of $p$ in $1,2,\dots,n$ are $p,2p,\dots$, and $2p>n$, so $p$ is the only one, and $p^2\nmid p$. Hence the exponent of $p$ in the prime factorisation of $n!$ is exactly $1$.

In a perfect power $m^k$ with $k\ge2$, every prime exponent is a multiple of $k$, so it is $0$ or at least $2$. Therefore $n!$ is not a perfect power. $\blacksquare$
